\subsection{非阿贝尔局部化}%
\label{sub:Non-abelian localization}

我们先从经典情形开始。设 \(X\) 是带有 \(\G_m\) 作用的光滑射影簇，设 \(F\) 是一个 \(\G_m\)-等变完美复形，并记
\[
    X^{\G_m}=\bigsqcup_\alpha Z_\alpha
\]
为不动点轨迹分解为连通分支的分解。在将表示环 \(R(\G_m)\) 局部化，使得下面的等变欧拉类成为可逆之后，我们有
\[
    \chi_{\G_m}(X,F)
    =
    \sum_\alpha
    \chi_{\G_m}\ab(
        Z_\alpha,
        F|_{Z_\alpha}\otimes e(N_{Z_\alpha/X})^{-1}
    ),
\]
其中
\[
    e(N_{Z_\alpha/X})
    =
    \lambda_{-1}(N_{Z_\alpha/X}^\vee)
    =
    \sum_{i\geq0}(-1)^i[\Lambda^iN_{Z_\alpha/X}^\vee].
\]

现在考虑叠的情形。设
\[
    \mf{X}
    =
    \mf{X}^{\mr{ss}}
    \sqcup
    \bigsqcup_\alpha\ev_1(S_\alpha)
\]
是一个 \(\Theta\)-分层。由于 \(\ev_1|_{S_\alpha}\) 是闭浸入，我们可以将 \(S_\alpha\) 与其在 \(\mf{X}\) 中的像等同。该层的中心是
\[
    Z_\alpha\coloneqq\sigma^{-1}(S_\alpha)\subset\ms{Grad}(\mf{X}),
\]
其中 \(\sigma\colon\ms{Grad}(\mf{X})\to\ms{Filt}(\mf{X})\) 由投影 \(\Theta\to B\G_m\) 诱导。

\begin{thm}[虚非阿贝尔局部化 {\cite[Theorem~5.1]{HLlocalization}}]
    设 \(\mf{X}\) 是带有 \(\Theta\)-分层
    \[
        \mf{X}
        =
        \mf{X}^{\mr{ss}}
        \sqcup
        \bigsqcup_\alpha S_\alpha,
    \]
    的拟光滑导出叠，其中我们将每个层与其在 \(\ev_1\) 下的像等同，并设 \(Z_\alpha\) 为其中心。按照 \(\G_m\)-权重将余切复形的限制分解为
    \[
        \L_{\mf{X}}|_{Z_\alpha}
        \simeq
        \L_\alpha^-
        \oplus
        \L_\alpha^0
        \oplus
        \L_\alpha^+,
    \]
    其中 \(\L_\alpha^-\)、\(\L_\alpha^0\) 与 \(\L_\alpha^+\) 的同调分别集中在负、零、正权重上。定义
    \[
        \E_\alpha
        \coloneqq
        \on{Sym}\ab(
            \L_\alpha^-
            \oplus
            (\L_\alpha^+)^\vee
        )
        \otimes
        \det(\L_\alpha^+)^\vee
        [-\on{rk}(\L_\alpha^+)]
        \in
        \ms{QCoh}(Z_\alpha).
    \]
    这里 \(\det(\L_\alpha^+)\) 与 \(\on{rk}(\L_\alpha^+)\) 分别表示完美复形 \(\L_\alpha^+\) 的行列式与虚秩。我们称一个向量空间复形是有限维的，如果其所有上同调群的直和是有限维的。
    设 \(F\in\ms{Perf}(\mf{X})\)。假设
    \[
        R\Gamma\ab(
            Z_\alpha,
            F|_{Z_\alpha}\otimes\E_\alpha
        )
    \]
    对每个 \(\alpha\) 都是有限维的，并且除有限多个 \(\alpha\) 外都是无环的。则 \(R\Gamma(\mf{X},F)\) 是有限维的当且仅当 \(R\Gamma(\mf{X}^{\mr{ss}},F)\) 是有限维的，并且
    \[
        \chi(\mf{X},F)
        =
        \chi(\mf{X}^{\mr{ss}},F)
        +
        \sum_\alpha
        \chi\ab(
            Z_\alpha,
            F|_{Z_\alpha}\otimes\E_\alpha
        ).
    \]
    这里拟光滑是指 \(\L_{\mf{X}}\) 是 Tor-幅度包含于 \([-1,1]\) 的完美复形。
\end{thm}


\subsection{几何不变量理论}%
\label{sub:Geometric invariant theory}

关于几何不变量理论，见~\cite{MFK1994}。下面的向量丛例子使用了 Harder--Narasimhan 滤过~\cite{HN1975}。

设 \(G\) 是约化群，\(X\) 是带有 \(G\)-等变线丛 \(\mc{L}\) 的射影概形。半稳定轨迹可以用 Hilbert-Mumford 判据来定义，它由所有满足以下条件的点组成：对所有的余特征标
\(\lambda \colon \G_m \to G\)，我们可以定义点
\[ x_0 = \lim_{t \to 0} \lambda(t) x \]
并定义量
\[ \mu^L(x, \lambda) \coloneqq \frac{-\on{wt}_{x_0}(\mc{L})}{\norm{\lambda}}. \]
一个点是半稳定的，如果对 \(G\) 的所有余特征标都有 \(\mu^L(x, \lambda) \leq 0\)。稳定轨迹由所有使得不等式严格成立且稳定化子为有限的点组成。

\begin{thm}[{\cite{MFK1994}}]
    稳定轨迹与半稳定轨迹在 \(X\) 中是开的。此外，\(X^{\on{ss}}/G\) 是好的商，而 \(X^s/G\) 是几何商。  
\end{thm}

\begin{exm}
    考虑亏格至少为 \(2\) 的曲线上秩为 \(r\)、次数为 \(d\) 的向量丛的模叠 \(\ms{Bun}_{r,d}\)。回顾半稳定轨迹定义为由斜率半稳定向量丛组成。事实上，这可以通过 GIT 得到。

    固定 \(C\) 上的一个丰富线丛 \(\mc{O}(1)\)。于是存在某个 \(n_0\)，使得对所有 \(n \geq n_0\) 及 \(E \in \ms{Bun}_{r,d}\)，\(E(n)\) 由整体截面生成。因此，我们可以考虑
    \[ H^0(C, E(n)) \twoheadrightarrow E(n). \]
    若令 \(V = \C^N\)，其中 \(N = h^0(E(n))\)，则存在满射
    \[ V \otimes \mc{O}_C(-n) \twoheadrightarrow E. \]
    因此，我们可以把 \(E\) 看作某个 Quot 概形的元素。在 Quot 概形上，我们可以考虑泛族
    \begin{equation*}
    \begin{tikzcd}
        V \otimes \mc{O}_C(-n) \ar[twoheadrightarrow]{r} \ar{d} & \mc{E} \\
        \ms{Quot}_{V \otimes \mc{O}_C(-n), p_E} \times C \ar{r}{p} & \ms{Quot}_{V \otimes \mc{O}_C(-n)}.
    \end{tikzcd}
    \end{equation*}
    现在定义线丛
    \[ \mc{L} \coloneqq \det R^0 p_* (\mc{E}(m)) \]
    其中 \(m \gg 0\)。我们也可以考虑 \(\mr{SL}(V)\) 在 Quot 概形上的作用。对某个余特征标 \(\lambda\)，我们可以分解
    \[ V = V_1 \oplus V_2 \]
    其权重分别为 \(s\) 和 \(N-s\)。
    然后考虑
    \[ V_1 \otimes \mc{O}_C(-n) \subset V \otimes \mc{O}_C(-n) \twoheadrightarrow E \]
    并考虑由它诱导的滤过
    \[ 0 \to F^1 E \to E \to E / F^1 E \to 0 \]
    于是我们看到
    \[ \lim_{t \to 0} \lambda(t) E = F^1 E \oplus E/F^1 E. \]
    接着我们计算限制 \(\mc{L}_{|\on{gr}_F E}\)，并得到空间 \(H^0(C, F^1 E(m))\) 与 \(H^0(C, E(m)/F^1 E(m))\)，因此 \(\mc{L}\) 的权重由
    \[ (N-s) p_{F^1 E}(m) - s p_{E/F^1 E}(m), \]
    给出，这重新得到了 GIT。
\end{exm}

现在考虑不稳定轨迹。设 \(E\) 是 \(C\) 上的向量丛。则存在唯一的滤过
\[ 0 \subset F^{\ell} E \subseteq \cdots \subseteq F^0 E = E \]
使得
\begin{enumerate}
    \item 所有 \(\on{gr}_F^i E\) 都是半稳定的；
    \item 对所有 \(i\)，有 \(\mu(\on{gr}_F^i E) > \mu(\on{gr}_F^{i-1} E)\)。 
\end{enumerate}
\begin{rmk}
    从 GIT 的角度看，我们实际上应当取 \(\mu\) 为约化希尔伯特多项式。
\end{rmk}

\subsection{\(\Theta\)-分层的构造}%
\label{sub:Construction of Theta-stratifications}

设 \(\mf{X}\) 是一个合理的代数叠。\footnote{这里，合理意味着局部有限表现且具有 qcqs 对角线之类的性质。}现在考虑 \(\A^1 \times C\) 上的某个族 \(\mc{E}\)。若限制到中心纤维，我们看到 \(\G_m\) 作用在 \(\on{gr}_F E\) 上，这意味着有一个映射 \(\G_m \to \Aut \on{gr}_F E\)，由此可以计算权重。

注意到态射 \(\sigma \colon \ms{Grad}(\mf{X}) \to \ms{Filt}(\mf{X})\) 是 \(\Theta\)-收缩，因此
\[ \pi_0(\ms{Filt}(\mf{X})) = \pi_0(\ms{Grad}(\mf{X})). \]

\begin{exm}
    设 \(\mf{X} = BG\)。则有
    \[\ms{Filt}(\mf{X}) = \bigsqcup_{[\lambda]} B P_{\lambda}, \]
    其中 \(P_{\lambda}\) 是对应于余特征标 \(\lambda\) 的抛物子群，并且我们总可以在单个 Weyl 室中选取代表元。此外，我们有
    \[ \ms{Grad}(\mf{X}) = \bigsqcup_{\lambda \in \mathbb{X}_*(T)/W} B L_{\lambda}, \]
    其中 \(L_{\lambda}\) 是对应于 \(\lambda\) 的 Levi 子群。
\end{exm}

\begin{defn}
    \(\mf{X}\) 中的一个 \textit{\(\Theta\)-层} 是连通分支的并 \(S \subseteq \ms{Filt}(\mf{X})\)，使得
    \[ \ev_1 \colon S \to \mf{X} \]
    是闭浸入。
\end{defn}

粗略地说，闭浸入是单射，因此 \(S\) 中的一个对象由 \(\mf{X}\) 中的一个对象以及一个“典范滤过”组成。

\begin{defn}
    一个 \textit{\(\Theta\)-分层} 是一族递增的开子叠
    \[ \mf{X}_{\leq c} \subseteq \mf{X}, \qquad c \in \Gamma, \]
    其中 \(\Gamma\) 是某个全序集，连同 \(\Theta\)-层
    \[ S_c \subseteq \ms{Filt}(\mf{X}_{\leq c}), \]
    它们被要求满足
    \[ \mf{X}_{\leq c} = \ev_1(S_c) \sqcup \bigcup_{c' < c} \mf{X}_{\leq c'}. \]
    特别地，这意味着
    \[ \mf{X} = \mf{X}^{\mr{ss}} \sqcup \bigsqcup_{c > 0} \ev_1(S_c). \]
\end{defn}

类似于 Hilbert-Mumford 判据，我们将利用数值不变量来构造 \(\Theta\)-分层。
\begin{defn}
    取值于全序集 \(\Gamma\) 的一个 \textit{数值不变量} \(\mu\) 由如下赋值组成：对任意 \(p \in \mf{X}(\C)\) 及任意非退化同态 \(\gamma \colon \G_m^q \to \Aut(p)\)，赋予一个尺度不变函数
    \[ \mu \colon \R^q \setminus \ab\{0\} \to \Gamma \]
    它们被要求满足
    \begin{enumerate}
        \item 该赋值在域扩张下不变。这一点不再详述，因为我们工作在 \(\C\) 上；
        \item 若有 \(\G_m^q \hookrightarrow \G_m^n \hookrightarrow \Aut(p)\)，则有分解 \(\R^q \setminus \ab\{0\} \to \R^n \setminus \{0\} \to \Gamma\)；
        \item 若有某个 \(\xi \colon S \to \mf{X}\) 及具有有限核的 \(\gamma \colon \G_{m, S}^q \to \Aut(\xi)\)，则函数 \(\mu_{\kappa(s)}\) 在 \(S\) 上局部常值，其中 \(s\) 取遍 \(S\) 的点。
    \end{enumerate}
\end{defn}

我们将给出一个朴素的构造版本，从数值不变量 \(\mu\) 出发。若对任意满足 \(f(1) = x\) 的非退化 \(f \in \ms{Filt}(\mf{X})\) 都有 \(\mu(f) \leq 0\)，则称 \(x \in \mf{X}\) 是半稳定的。若 \(x\) 不稳定，则定义其不稳定水平
\[ M^{\mu}(x) \coloneqq \sup \ab\{ \mu(f) \mid f \in \ms{Filt}(\mf{X}), \ev_1(f) = x\}. \]
若某个滤过 \(f\) 达到这个上确界，我们称之为“Harder-Narasimhan 滤过”。

\begin{defn}
    一个数值不变量 \(\mu\) 称为 \textit{标准的}，如果
    \begin{enumerate}
        \item 对 \(x \in \R^q \setminus \ab\{0\}\)，\(\mu(x)\) 与 \(\mu(-x)\) 不同时为正；
        \item 每个函数 \(\mu_{\gamma}\) 是严格拟凹的，即对所有 \(\delta \in [0,1]\) 及线性无关的 \(x_0, x_1 \in \R^q\)，若 \(\mu(x_0), \mu(x_1) \geq 0\)，则 
            \[\mu_{\gamma}(\delta x_0 + (1-\delta)x_1) \geq \max\ab\{\mu_{\gamma}(x_0), \mu_{\gamma}(x_1)\}.\] 若其中任一值严格为正，则要求该不等式严格成立。
    \end{enumerate}
\end{defn}

下一个条件是我们需要某种有理性。我们考虑 \(\mathbb{X}_*(T)_{\R}\)，目标是找到有理点。这引导我们得到条件 (R)：若 \(\mu_{\gamma}\) 对某个 \(\gamma\) 取到正值，则它在一个有理点处取到最大值。

\begin{thm}[{\cite{HLinstability}}]
    设 \(\mf{X}\) 是叠，\(\mu\) 是满足条件 (R) 的标准数值不变量。则 \(\mu\) 定义一个 \(\Theta\)-分层，当且仅当它满足以下条件：
    \begin{itemize}
        \item 对任意具有分式域 \(K\) 和剩余域 \(k\) 的 DVR \(R\)，有 \(\mu(f_K^{\mr{HN}}) \leq M^{\mu}(f_k)\)（HN-特殊化）；
        \item 对从有限型概形出发的任意映射 \(B \colon T \to \mf{X}\)，存在拟紧开子叠 \(\mf{X}' \subseteq \mf{X}\)，使得对任意有限型点 \(p \in T(k)\) 以及 \(B(p)\) 的一个滤过 \(f\)，若 \(\mu(f) > 0\)，则存在 \(B(p)\) 的另一个滤过 \(f'\)，满足 \(\mu(f') \geq \mu(f)\) 且 \(f'(0) \in \mf{X}'\)（HN-有界性）。
    \end{itemize}
\end{thm}

这里的要点是：标准性与 (R) 这两个条件结合 HN-有界性蕴含存在性，而 HN-特殊化蕴含唯一性。



